\Needspace{18\baselineskip}
\subsection{The maximal coefficient under positive lower mass}

For a Hermitian background $h$, the scalar in the preceding equation equals the subsolution threshold $b_h(f)$ from \cref{sec:individual-density-subsolutions}. One envelope-class solution bounds the coefficient of every subsolution by scalar comparison. Dominated continuity also identifies the limits of the normalized exponential solutions.

\begin{corollary}\label{cor:positive-lower-mass-exact-threshold}
Let $h$ be Hermitian with $v_-(h)>0$, and let $0\leq f\in L^1(\mathrm{d}V)$ be nonzero. There is $v\in\mathcal E_{\mathrm{env}}(X,h,0)$ such that
\[
 \mathrm{MA}^{\mathrm{ext}}_h(v)=b_h(f)f\,\mathrm{d}V.
\]
The number $b_h(f)>0$ is the unique positive scalar for which this equation has a solution in that envelope class.

Let $u_\lambda$ be the greatest exponential solution from \cref{thm:individual-exponential-subsolution}, and set $v_\lambda=u_\lambda-\sup_Xu_\lambda$. Each $u_\lambda$ belongs to $\mathcal E_{\mathrm{env}}(X,h,0)$. If $\lambda_j\downarrow0$ and $v_{\lambda_j}\to v$ in $L^1$, then the convergence holds in capacity, $v$ belongs to the same envelope class and satisfies the displayed equation.
\end{corollary}
\begin{proof}
By \cref{thm:general-measure-minimal}, there are $c>0$ and $w\in\mathcal E_{\mathrm{env}}(X,h,0)$ with $\mathrm{MA}^{\mathrm{ext}}_h(w)=cf\,\mathrm{d}V$. If $s\in\PSH(X,h)$ satisfies $\mathrm{MA}^{\mathrm{ext}}_h(s)\geq af\,\mathrm{d}V$ with $a>0$, then
\[
 \mathrm{MA}^{\mathrm{ext}}_h(w)
       \leq\frac ca\,\mathrm{MA}^{\mathrm{ext}}_h(s).
\]
Every $h$-plurisubharmonic $s$ satisfies $s\preceq0$. Thus the scalar assertion of \cref{cor:relative-extended-exponential-order}, with reference zero, gives $c/a\geq1$. Maximizing over subsolution coefficients yields $b_h(f)\leq c$. The reverse inequality holds because $w$ is itself a subsolution. The same argument identifies the scalar of every solution in this envelope class.

For $\lambda>0$, \cref{thm:general-measure-minimal} supplies an exponential solution $s_\lambda$ in the envelope class. The greatest-solution property gives $s_\lambda\leq u_\lambda$. Upward stability, \cref{prop:relative-elementary}, places $u_\lambda$ and its constant translate $v_\lambda$ in that class.

Write $c_\lambda=\exp(\lambda\sup_Xu_\lambda)$. By \cref{cor:individual-positive-threshold-exponential}, $c_\lambda\to b_h(f)>0$ and
\[
 \mathrm{MA}^{\mathrm{ext}}_h(v_\lambda)
       =c_\lambda\mathrm{e}^{\lambda v_\lambda}f\,\mathrm{d}V
       \longrightarrow b_h(f)f\,\mathrm{d}V
       \quad\text{in total variation}.
\]
Since $v_\lambda\leq0$ and $c_\lambda$ stays bounded for small $\lambda$, these measures are dominated by one finite non-pluripolar measure $C f\,\mathrm{d}V$. Apply \cref{thm:relative-fixed-measure-continuity} to the stated $L^1$ convergent sequence. It gives class membership, capacity convergence and the displayed equation for its limit.
\end{proof}

\Cref{thm:normalized-compactness} supplies such convergent subsequences. The final section gives examples with more than one normalized zero-exponent solution; accordingly, the corollary identifies cluster equations and does not assert convergence of the whole family.
